% ch8_b 第8章 §8.3–§8.4 (书页 345–353 / p360–p368)
% 实际覆盖范围说明：书页 345（p360）顶部为 8.2.2 节标题（非 8.3 节），故本块实际含
% 8.2.2、8.2.3、8.3、8.3.1、8.3.2、问题 8.3.1 与 8.4 节。
% 块界说明：书页 345 顶部为新节标题，无前块溢出内容，无 JOIN；
% 书页 353（p368）以一段正文自然收束全章，其后整页空白（p369 为第 9 章章首），无 TAIL。

\subsection{拓扑序的分类}

\begin{keypoints}
\item 理解拓扑序背后的数学框架是十分重要的。（正如理解群论——对称性破缺序背后的数学框架——十分重要一样。）
\item 对数学框架的理解将使我们能够对所有可能的拓扑序进行分类。
\end{keypoints}

如何标记和分类 FQH 液体中丰富的拓扑序，是一个长期悬而未决的问题。我们之所以能够对所有的晶体序进行分类，是因为我们知道晶体序由一个对称群描述。然而，我们对 FQH 液体中拓扑序的了解非常贫乏，拓扑序背后的数学结构尚不清楚。

尽管如此，对于一类 FQH 液体——阿贝尔 FQH 液体（Blok and Wen, 1990a,b; Read, 1990; Fröhlich and Kerler, 1991; Fröhlich and Studer, 1993）——我们还是找到了一种简单而统一的处理方法。Laughlin 态是最简单的阿贝尔 FQH 态，只含一种不可压缩流体组分。更一般的阿贝尔 FQH 态，其填充因子如 $\nu=2/5,3/7,\ldots$，含有数种不可压缩流体组分，并具有更复杂的拓扑序。阿贝尔 FQH 态中的拓扑序（或称跳舞图案（dancing patterns））同样可以用舞步来描述。跳舞图案可以用一个整数对称矩阵 $K$ 和一个整数电荷矢量 $\pmb{q}$ 来刻画。$\pmb{q}$ 的一个分量 $q_{i}$，是不可压缩流体第 $i$ 种组分中的粒子所携带的电荷（以 $e$ 为单位）。$K$ 的一个矩阵元 $K_{ij}$，是第 $i$ 种组分的粒子绕第 $j$ 种组分的粒子走一圈所迈的步数。在 FQH 态的 $(K,\pmb{q})$ 表征中，$\nu=1/m$ 的 Laughlin 态由 $K=m$、$\pmb{q}=1$ 描述，而 $\nu=2/5$ 的阿贝尔态由 $K=\begin{pmatrix}3&2\\2&3\end{pmatrix}$、$\pmb{q}=\begin{pmatrix}1\\1\end{pmatrix}$ 描述。

与拓扑序相关联的所有物理性质都可以用 $K$ 和 $\pmb{q}$ 定出。例如，填充因子就简单地由 $\nu=\pmb{q}^{\top}K^{-1}\pmb{q}$ 给出，而基态在亏格为 $g$ 的黎曼面上的简并度是 $\det(K)^{g}$。这一类 FQH 液体中所有的准粒子激发都具有阿贝尔统计，“阿贝尔 FQH 液体”由此得名。

上述对 FQH 液体的分类并不完全。并非每个 FQH 态都能用 $K$ 矩阵描述。1991 年，一类新的 FQH 态——非阿贝尔 FQH 态——被提出（Moore and Read, 1991; Wen, 1991b）。非阿贝尔 FQH 态含有具有非阿贝尔统计的准粒子。实验观察到的填充因子为 $\nu=5/2$ 的 FQH 态（Willett et al., 1987）很可能就是这样一个态（Haldane and Rezayi, 1988a,b; Greiter et al., 1991; Read and Green, 2000）。许多研究（Moore and Read, 1991; Blok and Wen, 1992; Iso et al., 1992; Cappelli et al., 1993; Wen et al., 1994）揭示了 FQH 态中的拓扑序与共形场论之间的联系。然而，对于非阿贝尔态中所有可能的拓扑序的完全分类，我们仍然相去甚远。

\subsection{边缘激发——测量拓扑序的一种实用方法}

\begin{keypoints}
\item FQH 态的边缘激发扮演着类似于 X 射线之于晶体的角色。我们可以利用边缘激发在实验上探测 FQH 态中的拓扑序。换言之，与基态简并度相比，边缘激发为拓扑序提供了更完全的定义。
\end{keypoints}

基态的拓扑简并只能部分地刻画拓扑序。不同的拓扑序有时会导致相同的基态简并度。这里的问题是，我们能否对拓扑序有更完全的刻画/测量。认识到拓扑序既不能用局域序参量、也不能用局域算符的长程关联来刻画之后，似乎很难再找到任何刻画拓扑序的方法。令人惊奇的是，FQH 态以一种出人意料的方式找到了出路。FQH 态体内的拓扑序可以用边缘激发来刻画/测量（Wen, 1992）。这种二维拓扑序编码在一维边缘态中的现象，与超弦理论和量子引力中的全息原理有若干相似之处（'t Hooft, 1993; Susskind, 1995）。\footnote{译注：原文为 holomorphic principle（全纯原理），疑为 holographic principle（全息原理）之误，此处按本意译出。}

作为不可压缩液体，FQH 液体的所有体内激发都有有限的能隙。然而，有限大小的 FQH 液体总是含有一维的无能隙边缘激发，这是 FQH 流体的另一个独特性质。边缘激发的结构极为丰富，这反映了体内丰富的拓扑序。不同的体内拓扑序导致不同的边缘激发结构。因此，我们可以通过研究边缘激发的结构来研究和测量体内的拓扑序。

正如我们在上一章所看到的，由于非平庸的体内拓扑序，（阿贝尔）FQH 液体边缘上的电子形成一类新的关联态——手征 Luttinger 液体（Wen, 1992）。手征 Luttinger 液体中的电子传播子获得一个反常指数：$\langle c^{\dagger}(t,x)c(0)\rangle\propto(x-vt)^{-g}$，$g\neq1$。（对费米液体，我们有 $g=1$。）在许多情形下，指数 $g$ 是一个拓扑量子数，不依赖于边缘的细节性质。因此，$g$ 是一个可用于刻画 FQH 液体中拓扑序的新量子数。许多实验组已经通过两条边缘之间隧穿电导的温度依赖关系成功地测量了指数 $g$（Milliken et al., 1995; Chang et al., 1996），理论预言其形式为 $\sigma\propto T^{2g-2}$（Wen, 1992）。这一实验证实了新型手征 Luttinger 液体的存在，并为在实验上研究 FQH 液体丰富的体内与边缘结构打开了大门。

%FIG{8.4}: {一维晶体经过杂质时，将在电压降中产生窄带噪声。}
%FIG{8.5}: {FQH 流体流经收缩区（constriction）时，由于准粒子的背散射而产生窄带噪声。}

非阿贝尔 FQH 液体的边缘态构成更为奇异的一维关联系统，这类系统迄今尚未被命名。人们发现，这些边缘态与 $1+1$ 维共形场论密切相关（Wen et al., 1994）。

我们知道，晶体序可以用 X 射线衍射实验来测量。下面我们想提出：FQH 液体中的拓扑序可以（原则上）通过边缘输运实验中的噪声谱来测量。让我们先考虑一个被驱动着通过杂质的一维晶体（见图 8.4(a)）。由于晶体序的存在，如果每个元胞只含一个带电粒子，杂质两侧的电压在频率 $f=I/e$ 处具有窄带噪声。更精确地说，噪声谱具有一条奇异性，即 $S(f)\sim A\delta(f-\frac{I}{e})$。如果每个元胞含有两个带电粒子（见图 8.4(b)），那么我们将在 $f=I/2e$ 处看到一条额外的窄带噪声，于是 $S(f)\sim B\delta(f-\frac{I}{2e})+A\delta(f-\frac{I}{e})$。在这个例子中我们看到，噪声谱使我们能够测量一维情形中的晶体序。类似的实验也可用于测量 FQH 液体中的拓扑序。让我们考虑一个带有窄收缩区的 FQH 样品（见图 8.5）。收缩区通过两条边缘之间的准粒子隧穿诱导背散射。背散射使收缩区两侧的电压出现噪声。在弱背散射极限下，噪声谱在某些频率处含有奇异性，这使我们能够测量 FQH 液体中的拓扑序。\footnote{此处给出的讨论仅适用于其边缘激发全部沿同一方向传播的 FQH 态。对于阿贝尔态，这要求 $K$ 的所有本征值都具有相同的符号。}更具体地，噪声谱中的奇异性具有如下形式（见式 (7.4.49)）
\begin{equation*}
S(f)\sim\sum_{a}C_{a}|f-f_{a}|^{\gamma_{a}}\tag{8.2.5}
\end{equation*}
奇异性的频率与指数 $(f_{a},\gamma_{a})$ 由拓扑序决定。对于由矩阵 $K$ 和电荷矢量 $\pmb{q}$ 表征的阿贝尔态，$(f_{a},\gamma_{a})$ 的允许取值为
\begin{equation*}
f_{a}=\frac{I}{e\nu}\pmb{q}^{\top}K^{-1}\pmb{l},\qquad\gamma_{a}=2\pmb{l}^{\top}K^{-1}\pmb{l}-1\tag{8.2.6}
\end{equation*}
其中 $\pmb{l}^{\top}=(l_{1},l_{2},\ldots)$ 是任意整数矢量，$\nu=\pmb{q}^{\top}K^{-1}\pmb{q}$ 是填充因子。噪声谱中的奇异性由两条边缘之间的准粒子隧穿引起。奇异性的频率 $f_{a}$ 由隧穿准粒子的电荷 $Q_{q}$ 决定，即 $f_{a}=\frac{I}{e}\frac{Q_{q}}{e\nu}$。指数 $\gamma_{a}$ 由隧穿准粒子的统计 $\theta_{q}$ 决定，即 $\gamma=2\frac{|\theta_{q}|}{\pi}-1$。于是，噪声谱测量的是允许存在的准粒子的电荷与统计，而这又决定了 FQH 态中的拓扑序。

\section{量子序}

\begin{keypoints}
\item 量子态一般包含一种新的序——量子序（quantum order）。量子序无法完全用破缺的对称性及相应的序参量来刻画。
\item 量子序描述多体基态中量子纠缠的图案。
\item 量子序的涨落可以产生无能隙规范玻色子和无能隙费米子。量子序保护这些激发的无能隙性，正如对称性保护对称性破缺态中的无能隙 Nambu--Goldstone 玻色子一样。
\item 拓扑序是量子序的一种特殊情形，其中所有激发都具有有限的能隙。
\end{keypoints}

按定义，拓扑序只描述有能隙量子态的内部序。这里我们做一次信念上的飞跃。我们将假设能隙并不重要，无能隙量子态同样可以包含无法用对称性和长程关联描述的序。我们将把量子基态中的非对称性破缺序称为量子序。

如果你相信这条思路，那么剩下要做的事情只有两件：证明量子序确实存在，并找到刻画量子序的数学描述（或符号）。我们将在 8.3.2 节和第 10 章中证明量子序确实存在。由于量子序不能用破缺的对称性和序参量来刻画，我们需要发展一个新的理论来描述量子序。目前，我们还没有一个能够描述所有可能量子序的完整理论。然而，在第 9 章中我们设法找到了一个数学对象——投影对称群（projective symmetry group, PSG）——它能够描述一大类量子序。

有人可能会问，为什么我们需要引入量子序这个新概念？它能有什么用？为了回答这类问题，我们想反过来问：为什么我们需要对称性破缺的概念？对称性破缺的描述有用吗？对称性破缺之所以有用，是因为它导致了对晶体序的分类（例如三维中 230 种不同的晶体），并且无需了解系统的细节就能定出低能激发的结构（例如固体中三支声子来自三个破缺的平移对称性）（Nambu, 1960; Goldstone, 1961）。量子序及其 PSG 描述在同一意义上是有用的：PSG 可以对具有相同对称性的不同量子态进行分类（Wen, 2002c），而量子序无需了解系统的细节就能定出低能激发的结构（Wen, 2002a,c; Wen and Zee, 2002）。对称性破缺序与量子序的主要差别在于：对称性破缺序产生并保护无能隙的 Nambu--Goldstone 模（Nambu, 1960; Goldstone, 1961），它们是标量玻色型激发；而量子序能够产生并保护无能隙规范玻色子和无能隙费米子。只要玻色子基态具有适当的量子序，费米子激发甚至可以在纯局域玻色模型中演生出来。

想象量子序的一种方式，是把量子序看作对多体基态中量子纠缠图案的描述。不同的纠缠图案给出不同的量子序。纠缠的涨落对应于量子有序态之上的集体激发。我们将看到，这些集体激发可以是规范玻色子和费米子。

拓扑序/量子序的概念在量子计算领域也十分有用。人们一直在设计各种不同的量子纠缠态来执行不同的计算任务。当量子比特的数目越来越大时，理解量子纠缠的图案就变得越来越困难。人们需要一个理论来刻画多量子比特系统中不同的量子纠缠。拓扑序/量子序理论（Wen, 1995, 2002c）正是这样一个理论。此外，Wen 和 Niu (1990) 所发现的拓扑有序态中稳健的拓扑简并，可以用于容错量子计算（Kitaev, 2003）。

下面我们将讨论量子相变与量子序之间的联系。然后，我们将利用自由费米系统中的量子相变来研究其中的量子序。

\subsection{量子相变与量子序}

\begin{keypoints}
\item 量子相变定义为基态能量作为哈密顿量中参数的函数所出现的奇异性。
\end{keypoints}

经典序可以通过经典相变来研究。经典相变以自由能密度 $f$ 中的奇异性为标志。自由能密度可以通过配分函数计算如下：
\begin{equation*}
f=-\frac{T\ln Z}{V_{\text{space}}},\qquad Z=\int\mathcal{D}\phi\,\mathrm{e}^{-\beta\int\mathrm{d}x\,h(\phi)}\tag{8.3.1}
\end{equation*}
其中 $h(\phi)$ 是经典系统的能量密度，$V_{\text{space}}$ 是空间的体积。

类似地，为了研究量子序，我们需要研究零温 $T=0$ 下的量子相变。这里，基态的能量密度扮演自由能密度的角色。基态能量密度中的奇异性标志着一次量子相变。基态能量密度与自由能密度之间的相似性，在下面基态能量密度的表达式中清楚可见：
\begin{equation*}
\rho_{E}=\mathrm{i}\frac{\ln Z}{V_{\text{space--time}}},\qquad Z=\int\mathcal{D}\phi\,\mathrm{e}^{\mathrm{i}\int\mathrm{d}x\,\mathrm{d}t\,\mathcal{L}(\phi)}\tag{8.3.2}
\end{equation*}
其中 $\mathcal{L}(\phi)$ 是量子系统的拉格朗日密度，$V_{\text{space--time}}$ 是时空的体积。比较式 (8.3.1) 与式 (8.3.2)，我们看到经典系统由一个正泛函的路径积分描述，而量子系统由一个复泛函的路径积分描述。一般而言，以复泛函路径积分的奇异性为标志的量子相变，可以比以正泛函路径积分的奇异性为标志的经典相变更为普遍。

%FIG{8.6}: {(a) 与 (b) 中的两组有向费米面代表两种不同的量子序。(a) 与 (b) 两种量子序之间两个可能的转变点由 (c) 与 (d) 中的费米面描述。}

\subsection{自由费米系统中的量子序与量子相变}

\begin{keypoints}
\item 自由费米系统包含不改变任何对称性的量子相变，这表明自由费米系统含有非平庸的量子序。
\item 自由费米系统中的不同量子序由费米面的拓扑来分类。
\end{keypoints}

让我们考虑一个只具有平移对称性以及来自费米子数守恒的 $U(1)$ 对称性的自由费米系统。哈密顿量具有如下形式：
\begin{equation*}
H=\sum_{\langle ij\rangle}\left(c_{i}^{\dagger}t_{ij}c_{j}+\text{h.c.}\right)
\end{equation*}
其中 $t_{ij}^{*}=t_{ji}$。基态由在每个负能量态上填充一个费米子得到。一般地，系统含有若干片费米面。

为了理解自由费米子基态中的量子序，我们注意到，当我们连续地改变 $t_{ij}$ 时，费米面的拓扑可以通过两种方式改变：一片费米面可以收缩为零（图 8.6(d)）；两片费米面可以连接起来（图 8.6(c)）。当 $d$ 维系统中的一片费米面即将消失时，基态能量密度具有如下形式：
\begin{equation*}
\rho_{E}=\int\frac{\mathrm{d}^{d}\pmb{k}}{(2\pi)^{d}}\left(\pmb{k}\cdot M\cdot\pmb{k}-\mu\right)\Theta\left(-\pmb{k}\cdot M\cdot\pmb{k}+\mu\right)+\cdots
\end{equation*}
其中"……"代表非奇异的贡献，对称矩阵 $M$ 是正（或负）定的。我们发现基态能量密度在 $\mu=0$ 处具有奇异性，即 $\rho_{E}=c\mu^{(2+d)/2}\Theta(\mu)+\cdots$，其中 $\Theta(x>0)=1$，$\Theta(x<0)=0$。当两片费米面即将连接时，奇异性仍由上式决定，只是此时 $M$ 既有负本征值又有正本征值。基态能量密度在 $d$ 为奇数时具有 $\rho_{E}=c\mu^{(2+d)/2}\Theta(\mu)+\cdots$ 形式的奇异性，在 $d$ 为偶数时具有 $\rho_{E}=c\mu^{(2+d)/2}\log|\mu|+\cdots$ 形式的奇异性。

%FIG{8.7}: {序的一种新分类。阴影框中的相可以用朗道的理论描述，其余的相则超越朗道的理论。}

基态能量密度在 $\mu=0$ 处的奇异性标志着一次量子相变。这类相变最早由 Lifshitz (1960) 研究。显然，相变两侧没有对称性的改变，也没有局域序参量可以表征相变两侧的相。这表明这两个态只能靠它们的量子序来区分。由于 $\mu=0$ 的位置恰好是费米面拓扑发生改变的地方，我们发现费米面的拓扑是一个刻画自由费米系统中量子序的"量子数"（见图 8.6）。拓扑的改变标志着改变量子序的连续量子相变。

\subsection*{问题 8.3.1}

考虑一个二维自旋-$1/2$ 自由电子系统。当我们改变化学势 $\mu$ 时，系统经历一次量子相变，如图 8.6(d) 所示。试求转变点 $\mu_{c}$ 附近自旋磁化率的奇异行为。

\section{序的一种新分类}

\begin{keypoints}
\item 量子序有许多类。FQH 态和自由费米系统只是众多类量子序中的两类。
\end{keypoints}

拓扑序/量子序的概念使我们能够对序进行一种新的分类，如图 8.7 所示。按照这一分类，量子序就是量子系统中的非对称性破缺序，而拓扑序就是具有有限能隙的量子序。

从上面讨论的 FQH 态和费米液体态可以看到，量子序可以分成若干不同的类。FQH 态和自由费米系统只提供了量子序的两个例子。在接下来几章中，我们将研究某些强关联系统中的量子序。我们将证明，这些量子序属于另外一个类，它与关联基态中弦网的凝聚密切相关。第 9 章将用投影构造（即从属玻色子方法）研究并分类这一类量子序。第 10 章将把第 9 章所研究的量子有序态与弦网凝聚态联系起来。
